Explain the basic idea of credibility theory, i.e., learning from the individual and
collective experience. How do the Bayesia/Frequentists approaches for net premia differ?

Credibility theory aims at finding risk-adequate premiums for different heterogeneous contracts in a
portfolio – given the observed claims history of the portfolio.

Let \( X_{i,t} \) represents the claim caused by contract \( i \) in the period \( t \), \( n_i \) is the sample size for the 
\( i \)-th contract.
then the frequentists approach for the net premia is

\( \overline{X}_i = \frac{1}{n_i} \sum_{k=1}^{n_i} X_{i,k} \simeq \mathbb{E}[X_{i, 1}] \)

Heterogeneity of contracts is captured by an unknown parameter for each contract that determines
the contract’s individual risk profile. The premium ansatz is chosen to be the expected claim size per
contract – given the knowledge on the contract’s risk profile.
Given the observed claims history of the overall portfolio, the premiums are estimated.
We choose the Bayes estimator (under SE) for estimating the premium.


Define the heterogeneous risk model. Explain that the relevant
information is characterized by the individual claims experience of the relevant
contract alone.

The question also included: The net risk premium is given by a conditional expectation given the true type and motivated by a conditional LLN. 
Discuss this.

The Heterogenous risk model is with the Prob. space \((\Omega, \mathcal{F}, \mathbb{P})\):
1. The \(i\)th risk is described through the random variables \((\theta_i, (X_{i,t})_{t=1,...,n})\). Here, \(\theta_i\) denotes the real-valued parameter of heterogeneity and \((X_{i,t})_{t=1,...,n}\) represent the claim sizes caused by risk \(i\) in over time.
2. \((\theta_i, (X_{i,t})_{t=1,...,n}), i = 1, ..., r\), are i.i.d. random elements.
3. It holds: \(X_{i,t} \in L^2 (\Omega, \mathcal{F}, \mathbb{P})\) for all \(i, t\).
4. Given \(\theta_i\), the sequence \((X_{i,t})_{t=1,...,n}\) is i.i.d. with conditional distribution function denoted by \(F(\cdot | \theta_i)\).

Which allows for the following Ansatz, that the net risk premium given \(\theta_i\) is:
\[
\mu(\theta_i) := E[X_{i,t} | \theta_i] = \int_{\mathbb{R_+}} x \, dF(x | \theta_i).
\]
Remark: Motivated by the law of large numbers: Given \(\theta_i\), it holds that
\[
\bar{X}_i := \frac{1}{n_i} \sum_{t=1}^{n_i} X_{i,t} \xrightarrow[n_i \to \infty]{} \mu(\theta_i) \quad \mathbb{P}\text{-a.s.},
\]
i.e., \(\mu(\theta_i)\) is an approximation of the average claim size \(\bar{X}_i\) of risk \(i\) given \(\theta_i\).
Estimate the net risk premium \( \mu(\theta_i) \) given the claims history of the overall portfolio
\( \mathfrak{X}_j = (X_{j,1}, \dots, X_{j,nj})^\top \), \(j = 1, \dots, r\).

Idea: Estimation through minimization of the \(L^2\) distance:
\(
\rho(\hat{\mu}) := E[(\mu(\theta_i) - \hat{\mu})^2].
\)
Overall, we thus obtain the following optimization problem
\(
\text{minimize } \rho(\hat{\mu}) = E[(\mu(\theta_i) - \hat{\mu})^2] \text{ over all } \hat{\mu} \in L^2 (\Omega, \mathcal{F}_0, \mathbb{P}).
\)
Note that this is just the Bayes estimator with SE as Loss function in disguise. As a consequence we get

Theorem:
The optimization problem possesses a solution \(\hat{\mu}^{\ast} \in L^2 (\Omega, \mathcal{F}_0, \mathbb{P})\) namely the conditional expectation
\(
\hat{\mu}^{\ast} := E[\mu(\theta) | \mathcal{F}_0] = E[\mu(\theta_i) | X_i].
\)
The solution is \(\mathbb{P}\)-a.s. unique. The minimal squared error (minimal Bayes risk) is given by
\(
\rho(\hat{\mu}^{\ast}) = E[\text{Var}[\mu(\theta_i) | X_i]],
\)
where \(\text{Var}[\mu(\theta_i) | X_i] := E[(\mu(\theta_i) - E[\mu(\theta_i) | X_i])^2 | X_i]\) is the conditional variance of \(\mu(\theta_i)\) given the data \(X_i\).

The Theorem shows that the estimator \(\hat{\mu}^{\ast}\) for \(\mu(\theta_i)\) given the claims experience \( X_1, \dots, X_r \) of all risks only
depends on the individual claims experience \( X_i = (X_{i,1}, \dots, X_{i, n_i})^\top \) of risk \( i \).
To explicitly calculate the Bayes estimator, we need to know the (posterior) distribution of \( \thea_i|\mathcal{X}_i \).

Discuss Bayesian estimators in the context of credibility theory with the Poisson sampling model.

We assume that the parameter of heterogeneity \(\theta\) is \(\Gamma(\alpha, \beta)\)-distributed (prior distribution) and that given \(\theta\), 
\(X_t\), \(t = 1, ..., n\), are independent and identically \(\text{Poi}(\theta)\)-distributed (we thus consider a Poisson sampling model).

From the previous sections, we know that in this situation, \(X \sim \text{NB}(\alpha, \frac{\beta}{1+\beta})\) and it holds that
\[
\mu(\theta) = E[X_1 | \theta] = \theta.
\]
To obtain our estimator \(\mu^{\ast} = E[\mu(\theta) | X] = E[\theta | X]\) explicitly, we first calculate the density of the (posterior) distribution 
\(f_{\theta | X}\).

Since \(f_{\theta}(y) = \frac{\beta^{\alpha}}{\Gamma(\alpha)} y^{\alpha-1} e^{-\beta y} \mathbb{I}_{(0, \infty)}(y)\), the posterior is of the form
\[
f_{\theta | X}(y, X) = \frac{1}{P[X=x]} P[X_1 = x_1 | \theta] \cdots P[X_n = x_n | \theta] f_{\theta}(y)
\]
\[
= C_1(x) y^{\alpha-1} e^{-\beta y} \prod_{t=1}^{n} \frac{y^{x_t}}{x_t!} e^{-y}
\]
\[
= C_2(x) y^{\alpha + \sum_{t=1}^{n} x_t - 1} e^{-(\beta + n)y}
\]
for constants \(C_1(x)\), \(C_2(x)\).

This is of the form of a gamma distribution, i.e.,
\[
\theta | X \sim \Gamma \left( \alpha + \sum_{t=1}^{n} X_t, \beta + n \right).
\]
Using the formula for the expected value of the gamma distribution, we obtain
\[
\hat{\mu}^{\ast} = E[\theta | X] = \frac{\alpha + \sum_{t=1}^{n} X_t}{\beta + n}.
\]
Since \(E[\theta] = \frac{\alpha}{\beta}\), the Bayes estimator possesses the alternative representation
\[
\hat{\mu}^{\ast} = (1 - w) E[\theta] + w \bar{X}
\]
with weight \(w := \frac{n}{\beta + n} \in (0, 1)\) and \(\bar{X} = \frac{1}{n} \sum_{t=1}^{n} X_t\).
The minimal squared error according to our result for the optimal solutionis given by
\[
\rho(\mu^{\ast}) = E[\text{Var}[\mu(\theta) | X]] = E[\text{Var}[\theta | X]]
\]
\[
= E \left[ \frac{\alpha + \sum_{t=1}^{n} X_t}{(\beta + n)^2} \right] = \frac{\alpha + n E[X_1]}{(\beta + n)^2}
\]
\[
= \frac{\alpha}{\beta} \frac{1}{\beta + n},
\]
where the last equality follows from \(E[X_1] = E[E[X_1 | \theta]] = E[\theta] = \frac{\alpha}{\beta}\).
As one would expect - the minimal squared error, i.e., the estimation error, decreases with increasing number of observation periods \(n\).


Explain the reason for the best linear estimator in credibility theory and its definition.
Explain in terms of L2 -projection operators the difference to the usual estimator.

In general, calculating the estimator \( \hat{\mu}^{\ast} = E[\mu(\theta_i) | X_i] \) is only feasible for some distributions,
A workaround is restricting the set of allowed functions for \( \hat{\mu}^{\ast} \), namely the subset of linear estimators:
\(\mathcal{L} := \{ \hat{\mu} | \hat{\mu} = a_0 + \sum_{i=1}^{r} \sum_{t=1}^{n_i} a_{i,t} X_{i,t} \text{ with } a_0, a_{i,t} \in \mathbb{R} \}\)

The reformulated goal is to obtain the best linear estimator of \(X\) given \(Y_1, ..., Y_m\) using the \(L^2\) distance, 
i.e., we want to find an estimator
\[
\hat{Z}_L^{\ast} \in \mathcal{C}' := \{ \hat{Z} \mid \hat{Z} = a_0 + a^T Y \text{ with } a_0 \in \mathbb{R}, a \in \mathbb{R}^m \}
\]
that minimizes \(E[(X - \hat{Z})^2]\) over all \(\hat{Z} \in \mathcal{C}'\).

Describe the general characterization of a best linear estimator. Note that this
involves only moments up to order 2.

Let \((a_0, a)\) be an arbitrary solution to the system of linear equations
\[
a_0 = E[X] - a^T E[Y], \quad \Sigma_{X,Y} = \Sigma_Y a
\]
and \(\hat{Z}_L := a_0 + a^T Y\). Then, for all \(\hat{Z} \in \mathcal{C}'\), it holds that
\[
E[(X - \hat{Z})^2] \ge E[(X - \hat{Z}_L)^2] = \text{Var}[X] - a^T \Sigma_Y a,
\]
i.e., each pair \((a_0, a)\) corresponds to a best linear estimator \(\hat{Z}_L\). 
Conversely, the equations (10) also constitute a necessary condition for a best linear estimator.

An estimator \(\hat{Z}_L^{\ast}\) from a) fulfills the equations
\(
E[X] = E[\hat{Z}_L^{\ast}] \quad \text{and} \quad \text{Cov}[X, Y_i] = \text{Cov}[\hat{Z}_L^{\ast}, Y_i] \quad (i = 1, ..., m).
\)
If \(\Sigma_Y\) is invertible, there exists a unique minimizer \(\hat{Z}_L^{\ast}\) of \(E[(X - \hat{Z})^2]\) in the class \(\mathcal{C}'\), namely
\(
\hat{Z}_L^{\ast} = E[X] + \Sigma_{X,Y} \Sigma_Y^{-1} (Y - E[Y]).
\)
The corresponding minimal squared error is given by
\(
E[(X - \hat{Z}_L^{\ast})^2] = \text{Var}[X] - \Sigma_{X,Y} \Sigma_Y^{-1} \Sigma_{X,Y}^T = \text{Var}[X] - \text{Var}[\hat{Z}_L^{\ast}].
\)

State the assumptions and results of the Bühlmann model
and the Bühlmann-Straub model

The Bühlmann Model is defined through the following assumptions:

1. The \(i\)th contract is described through \((\theta_i, (X_{i,t})_{t=1,...,n})\). Here, \(\theta_i\) denotes the 
parameter of heterogeneity and the random variables \((X_{i,t})_{t=1,...,n}\) represent claim sizes of the \(i\)th contract over time.
2. The pairs \((\theta_i, (X_{i,t})_{t=1,...,n})\), \(i = 1, ..., r\), are independent.
3. The sequence \((\theta_i)_{i=1,...,r}\) is i.i.d.
4. It holds \(\text{Var}[X_{i,t}] < \infty\) for all \(i\) and \(t\).
5. Given \(\theta_i\), the random variables \(X_{i,t}\), \(t \ge 1\), are independent and their expected values and variances are functions of 
\(\theta_i\):
\(
\mu(\theta_i) = E[X_{i,t} | \theta_i] \quad \text{and} \quad v(\theta_i) = \text{Var}[X_{i,t} | \theta_i].
\)

Technical Remark:
Due to 2.), \((\mu(\theta_i))_{i=1,...,r},\ (v(\theta_i))_{i=1,...,r}\) are i.i.d. sequences
We will write:
\(
\mu = E[\mu(\theta_i)], \quad \lambda = \text{Var}[\mu(\theta_i)] \quad \text{and} \quad \varphi = E[v(\theta_i)].
\)

The central innovation over the Heterogenous model is that the individuals/risk profiles \( X_i \) can have different distributions
(which does not vary over time though).
One can then derive from the covariance structure (\( \text{Cov}[X_{i,t}, X_{i, s}] = \lambda + \phi \) for \( t = s \) or else \( \lambda \),
that under the assumption \(\varphi > 0\), a best linear estimator \(\hat{\mu}_L = a_0 + a^T \text{ of } \mu(\theta_i)\) given the data 
\(\bar{X}_1, ..., \bar{X}_r\) exists. It is \(\mathbb{P}\)-a.s. unique and given by
\(
\hat{\mu}_L^{\ast} = (1 - w) \mu + w \bar{X}_i
\)
where
\[
w = \frac{n_i \lambda}{\varphi + n_i \lambda} \in (0, 1) \quad \text{and} \quad \bar{X}_i = \frac{1}{n_i} \sum_{t=1}^{n_i} X_{i,t}.
\]. \(\hat{\mu}_L^{\ast}\) is called the credibility estimator and \( w \) the credibility weight, 
the higher it is the stronger the emphasis on individual claims, 
in fact with \( w \to 1, n_i \to \infty \) this becomes the frequentist approach.
The minimal squared error is given by
\(
\rho(\hat{\mu}_L^{\ast}) = (1 - w) \lambda.
\).

If we change the property 5 in the Bühlmann Model we get the Bühlmann-Straub Model, in particular
5'. Given \(\theta_i\), \(X_{i,t}\), \(t \ge 1\), are independent; their expectation and variances are given as functions of \(\theta_i\):
\[
\mu(\theta_i) = E[X_{i,t} | \theta_i] \quad \text{and} \quad \text{Var}[X_{i,t} | \theta_i] = \frac{v(\theta_i)}{p_{i,t}},
\]
where \(p_{i,t}\) are known positive weights and let \( p_i = \sum_{t=1}^{n_i} p_{i,t} \).
The main innovation here is heterogeneity over time for an individual/risk profile.

The best linear estimator is:
If the covariance matrix \(\Sigma_X\), \(i \in \{1, ..., r\}\) is invertible, it is
\(
\hat{\mu}_L^{\ast} = (1 - w) \mu + w \bar{X}_i
\)
where
\(
w = \frac{\lambda p_i}{\varphi + \lambda p_i} \quad \text{and} \quad \bar{X}_i = \frac{1}{p_i} \sum_{t=1}^{n_i} p_{i,t} X_{i,t}
\)
the \(\mathbb{P}\)-a.s. unique best linear estimator of \(\mu(\theta_i)\) given the data \(\bar{X}_1, ..., \bar{X}_r\). 
The minimal squared error is given by
\(
\rho(\hat{\mu}_L^{\ast}) = (1 - w) \lambda.
\)


\iffalse
Due to 2.), claim histories belonging to different risks are independent.
Claim sizes belonging to a single risk \(X_{i,1}, ..., X_{i,n_i}\), however, can be correlated:

Let \(\mu(\theta_i) := E[X_{i,1} | \theta_i]\). Then, because \(E[X_{i,s} | \theta_i] = E[X_i | \theta_i] = \mu(\theta_i)\) and \(E[X_{i,s}] = E[E[X_{i,s} | \theta_i]] = E[\mu(\theta_i)] = E[X_{i,t}]\) we obtain from 4.) that
\[
\text{Cov}[X_{i,s}, X_{i,t}] = E[E[(X_{i,s} - E[X_{i,s}])(X_{i,t} - E[X_{i,t}]) | \theta_i]]
\]
\[
= E[(E[X_{i,s} | \theta_i] - E[X_{i,s}])(E[X_{i,t} | \theta_i] - E[X_{i,t}])]
\]
\[
= \text{Var}[\mu(\theta_i)],
\]
i.e., \(\text{Cov}[X_{i,s}, X_{i,t}] > 0\) for \(\text{Var}[\mu(\theta_i)] > 0\).

Claim sizes \(X_{i,t}\), \(t = 1, ..., n_i\), \(i = 1, ..., r\), are identically distributed since for all \(x \in \mathbb{R}\):
\[
F_{X_{i,t}}(x) = P[X_{i,t} \le x] = E[P[X_{i,t} \le x | \theta_i]] \overset{4.)}{=} E[F(x | \theta_i)] \overset{2.)}{=} E[F(x | \theta_1)].
\].
\fi



